\subsection{Transformer and ClimateBERT Models for ESG ABSA}

\subsubsection{Overview: BERT and domain-adapted NLP foundations}

Foundation and rationale. BERT and its descendants provide deep bidirectional representations that capture context across long textual spans, enabling fine-grained ABSA where sentence- or clause-level sentiment is tied to explicit ESG aspects. Domain-adapted variants (e.g., FinBERT, ClimateBERT) initialize from general-purpose encoders but continue pretraining on large domain-specific corpora to inject specialized vocabulary, discourse patterns, and regulatory semantics into representations. This paradigm underpins effective ESG ABSA in multilingual, regulatory-grounded settings by aligning model priors with domain language and task structure \cite{10.21608/msaeng.2023.291899,10.1108/ijaim-08-2023-0206,10.47473/2020rmm0103,10.18653/v1/2023.acl-short.91}
Practical enablement for ESG ABSA. With Transformers as a backbone, one can (i) perform aspect-term extraction and (ii) predict aspect-level sentiment conditioned on regulatory context, all while supporting cross-language transfer when combined with multilingual encoders (e.g., mBERT, XLM-R). The literature shows domain-adapted pretraining yields tangible gains over vanilla BERT in ESG/climate tasks, improving both classification accuracy and interpretability when coupled with task-specific heads or adapters \cite{10.21608/msaeng.2023.291899,10.18653/v1/2023.acl-short.91,10.1051/e3sconf/202343602004,10.48550/arxiv.2510.22964
}
\subsubsection{ClimateBERT: training data, task coverage, and known strengths}

Training data foundations. ClimateBERT is a DistilRoBERTa-based model pre-trained on climate-related text drawn from climate-news articles, climate policies, and climate-relevant reports, with a substantial corpus amounting to millions of climate-focused tokens. This domain-specific pretraining yields representations attuned to climate discourse, terminology, and argumentative patterns encountered in ESG communications and policy documents \cite{10.1145/3630106.3658966,10.48550/arxiv.2108.01415,10.3390/su15053919} . ClimateBERT’s design emphasizes climate-risk signals, commitments, and disclosure language that appear in corporate reporting and regulatory materials.
Task coverage and capabilities. Fine-tuning ClimateBERT has been demonstrated across multiple climate-related classification tasks, including climate-detection in paragraphs, sentiment regarding climate risks and opportunities, detection of climate commitments and actions, and classification of climate-specific disclosure categories (e.g., net-zero, reduction targets, renewable energy mentions). Extensions include net-zero target detectors (ClimateBERT-NetZero) that identify explicit net-zero or reduction targets and can be integrated with QA modules for targeted evidence extraction. These capabilities collectively support climate-focused ABSA and regulatory-aligned signal extraction in ESG contexts \cite{10.18653/v1/2023.emnlp-main.975,10.1371/journal.pone.0288052,10.48550/arxiv.2511.00024}
Strengths observed. Empirical studies consistently show ClimateBERT outperforms general-purpose models (e.g., vanilla BERT, uncased variants) on climate-disclosure tasks, particularly in classification accuracy, domain-relevant linguistic nuance, and efficiency of adaptation (domain-adaptive pretraining is cost-effective and yields robust gains) \cite{10.3390/su15053919,10.1371/journal.pone.0288052,10.48550/arxiv.2511.00024,10.48550/arxiv.2108.01415,10.1145/3630106.3658966}. When paired with external validators (e.g., ClimateBERT-NetZero), there is enhanced capability to detect and quantify commitments and targets across climate narratives.
Important caveat for ABSA. ClimateBERT is not a direct tone classifier; rather, it serves as a validation and grounding lens for climate-related signals. It excels at detecting climate-relevant content and targets but does not inherently provide a complete tone or aspect-sentiment mapping without additional ABSA heads or ontology-based integration. This distinction is crucial when interpreting outputs in a broader ESG ABSA pipeline \cite{10.18653/v1/2023.emnlp-main.975,10.1371/journal.pone.0288052,10.48550/arxiv.2511.00024} .
\subsubsection{Other domain-adapted models: FinBERT and ESG-BERT}

FinBERT and ESG-specific variants. FinBERT is a BERT-based model pre-trained on financial-domain text, optimized for sentiment analysis within financial narratives. Its strength lies in handling financial-lexicon, investor-oriented expressions, and context where ESG signals intersect with financial analysis. ESG-BERT variants extend this approach to ESG-specific corpora, aiming to capture domain-specific vocabulary and discourse patterns in environmental, social, and governance reporting. Empirical work indicates domain-adapted models provide superior sentiment discrimination in ESG-related content compared to general-purpose baselines, especially when restricted labeled data exist \cite{10.21608/msaeng.2023.291899,10.3390/math12193119}
Practical utility for ESG ABSA. FinBERT- and ESG-BERT-inspired architectures can be integrated into ABSA pipelines as specialized componentry: (i) an aspect-term extractor trained or fine-tuned on ESG taxonomies, (ii) a sentiment classifier leveraging finance- or ESG-tuned encodings, and (iii) a regulatory anchor linkage module to map outputs to GRI/SASB/TCFD concepts. The domain alignment improves performance on aspect-specific sentiment, particularly for financially material ESG topics, compared with vanilla BERT baselines \cite{10.21608/msaeng.2023.291899,10.3390/math12193119}
Limitations to consider. Domain-adapted models are data-hungry and may overfit to their training corpora if not carefully regularized or if cross-domain transfer is attempted. Additionally, while they improve Sentiment classification in their target domains, they require careful alignment with ESG-specific taxonomies and regulatory ontologies to ensure interpretability and auditability in bilingual Indonesian/English contexts \cite{10.1177/09726225241237304,10.1111/ejed.70225,10.1145/3503044}
\subsubsection{Limitations of ClimateBERT for broader ESG analysis}

Narrow domain specialization. ClimateBERT is optimized for climate-related content and the 14 or so Climate/T CFD categories it has been tuned to; its coverage of broader ESG domains (social governance, labor relations, diversity, governance processes beyond climate) is comparatively limited. For comprehensive ESG ABSA, reliance on ClimateBERT alone may miss non-climate environmental topics (e.g., biodiversity, supply chain labor standards, governance risk) unless augmented by additional domain-adapted or multi-task models \cite{10.3390/su15053919,10.1371/journal.pone.0288052,10.48550/arxiv.2511.00024}
Transfer and cross-language constraints. ClimateBERT’s in-domain pretraining is climate-centric and primarily English-language data or translation-augmented corpora; applying ClimateBERT directly to Indonesian-language sustainability disclosures may degrade performance unless addressed with multilingual adapters, cross-lingual fine-tuning, or explicit Indonesian-language climate datasets. Cross-language alignment is essential for bilingual ESG ABSA but is not guaranteed by ClimateBERT’s native pretraining alone \cite{10.1016/j.inffus.2025.103073,10.1111/eufm.12509,10.48550/arxiv.2511.00024}.
Limitations as a validation lens, not a sole classifier. The literature clarifies that ClimateBERT is not a turnkey tone classifier; it functions best as a validation or targeting tool for climate-specific content and commitments. For ABSA pipelines requiring aspect-level sentiment, ClimateBERT should be integrated with ontology-guided ABSA components and domain-specific taxonomies to yield interpretable, regulator-friendly outputs rather than stand-alone sentiment scores \cite{10.18653/v1/2023.emnlp-main.975,10.1371/journal.pone.0288052}
Considerations for governance and regulatory signaling. ClimateBERT’s climate-centric focus may underrepresent governance or social aspects critical to Indonesian OJK/IDX reporting and SASB-like requirements. Consequently, ABSA pipelines must incorporate broader ESG ontologies and cross-framework mappings (GRI/SASB/TCFD) to ensure that non-climate disclosures are properly captured and interpreted \cite{10.1177/09726225241237304,10.1111/ejed.70225,10.1145/3503044}.
\subsubsection{Practical recommendations for ESG ABSA pipelines incorporating transformer and ClimateBERT systems}

Hybrid architecture strategy. Combine climate-focused validators (ClimateBERT or ClimateBERT-NetZero) with broader ESG ABSA modules leveraging domain-adapted models (FinBERT/ESG-BERT) and a strong ontology-based information extraction layer. This hybrid approach allows climate commitments/targets to be cross-validated while maintaining coverage of broader ESG topics through ontology paths and ABSA heads.
Multilingual and cross-framework integration. For Indonesian/English bilingual disclosures, deploy multilingual transformers (e.g., XLM-R, mBERT) with domain adapters and cross-language alignment mechanisms to maintain consistent aspect labeling and climate signals across languages. Ontology-based mappings should connect climate signals to TCFD and SASB indicators and align Indonesian terms with international ESG vocabulary \cite{10.1016/j.inffus.2025.103073,10.1111/eufm.12509,10.48550/arxiv.2511.00024,10.1002/csr.70116}
Explainability and auditability. Leverage ClimateBERT’s outputs as a validation lens and fuse them with ontology-based explanations (e.g., path-based reasoning from climate mentions to net-zero targets) and lexical triggers to satisfy regulator demands for interpretability. Use ABSA-specific explainability tools (SHAP/LIME) alongside ontology-path visualizations to provide transparent audit trails \cite{10.1111/1475-679x.12575,10.3390/su11041093}
Evaluation strategy. Employ cross-domain benchmarks that test climate signal detection, net-zero/target extraction, and broader ESG aspect sentiment; perform ablation studies to isolate the contributions of ClimateBERT vs domain-adapted ESG models; and validate outputs against regulatory lexicons and target KPIs to ensure policy alignment and practical utility \cite{10.18653/v1/2023.emnlp-main.975,10.1371/journal.pone.0288052,10.48550/arxiv.2511.00024}.
\subsubsection{Summary: positioning ClimateBERT within a robust ESG ABSA ecosystem}

ClimateBERT serves as a strong, validated lens for climate-related content within ESG ABSA pipelines, offering domain-specific detection of environmental claims, climate risk discourse, and net-zero/renewable indicators. However, for holistic ESG ABSA—encompassing social and governance dimensions as well as cross-language Indonesian/English disclosures—ClimateBERT should operate as a complementary validator rather than a standalone solution. Integrating ClimateBERT with transformer-based general ABSA models, ontology-driven information extraction, and knowledge-graph-backed explainability yields a scalable, regulator-aligned framework capable of nuanced aspect-level sentiment analysis across the full ESG spectrum, while preserving cross-language robustness and interpretability \cite{10.18653/v1/2023.emnlp-main.975,10.1371/journal.pone.0288052,10.48550/arxiv.2511.00024,10.1145/3503044,10.1111/1475-679x.12575}